\section{Introduction}

Attention-deficit/hyperactivity disorder (ADHD) is characterized by persistent inattention and/or hyperactivity-impulsivity associated with functional impairment. Difficulties in executive control, persistence after error, frustration tolerance, and emotion regulation are also common in affected children, although emotion dysregulation is not itself a defining diagnostic criterion \parencite{graziano2016dysregulation,wolraich2019guideline}.

Digital interventions provide a useful experimental substrate because task contingencies, challenge intensity, sensory cues, and response timing can be standardized and logged. Randomized trials and reviews show that game-based and computerized interventions can improve selected trained or near-transfer outcomes in pediatric ADHD \parencite{kollins2020stars,oh2024dtx,lin2025seriousgames,yan2026ceft}. EndeavorRx, evaluated in the STARS-ADHD programme, is an important example that digital game-like delivery can support a regulated therapeutic product, but improvement on a targeted attention measure does not by itself establish broad transfer to everyday emotional or functional outcomes \parencite{kollins2020stars}.

The central problem addressed here is therefore not whether a child can become better at a game. It is whether a digital intervention can repeatedly train a measurable transition from challenge back to goal-directed behavior, and whether change in that transition is associated with independent outcomes outside the trained task.

ARC-ADHD treats regulation as a control problem. The child encounters a bounded challenge, interrupts an ineffective response, identifies relevant task or body information, selects a strategy, regulates, and resumes the original goal. Error, frustration, or temporary destabilization are not defined as intervention failure. The key event is successful recovery.

The present paper has two purposes. First, it formalizes the core ARC-ADHD mechanism and its observable quantities as an implementable computational research architecture. Second, it defines a staged evaluation programme that can falsify the central claim before a confirmatory efficacy trial is attempted. It is intentionally not presented as a completed randomized-trial protocol. A future confirmatory protocol should be written after the intervention version, primary endpoint, recovery parameters, active control, and analysis rules have been frozen, and should follow the then-current SPIRIT guidance \parencite{chan2025spirit}.

The computational contribution is fourfold. First, the intervention is represented as a versioned closed-loop state sequence with explicit challenge eligibility, regulation transitions, recovery monitoring, and return-to-goal states. Second, recovery is represented as a reconstructable event-level time-to-event quantity with target-rule-matched normalization, persistence, and right-censoring rather than a single postchallenge score. Third, the measurement context is locked while recovery is assessed so adaptive changes, safety interventions, telemetry loss, ordinary censoring, and scientifically invalid events remain distinguishable. Fourth, the framework couples this measurement contract to deterministic controller replay, versioned telemetry, and explicit BYOD qualification so the delivered intervention and derived recovery records can be reconstructed from the study archive.

Four technology-specific lines have been removed from the principal causal claim. Facial-behavior feedback, bilateral auditory stimulation, embodied/augmented-reality context transfer, and wearable EEG measurement are developed as separate companion papers. This separation reduces scope creep and prevents a null or positive result for one optional technology from being conflated with evidence about the core regulation mechanism.
